% =========================================================================
\section{Señales y amplificadores}\label{sec:amplificadores}
% =========================================================================
% -------------------------------------------------------------------------
\subsection{Valores característicos de una señal}\label{sec:senales}
% -------------------------------------------------------------------------
En el laboratorio una señal es una tensión $y(t)$ que el osciloscopio representa frente al tiempo.
Para compararla con los cálculos se resume en unos pocos números~\cite{TEORIA-EA}.

\begin{definicion}[Valores de una señal periódica]\label{def:valores}
Sea $y(t)$ periódica de periodo $T=1/f$. Se definen el valor de pico, el valor pico a pico, el
valor medio y el valor eficaz (\emph{root mean square}):
\begin{gather}
  y_{p}=\max_{t}\abs{y(t)},\qquad
  y_{pp}=\max_{t}y(t)-\min_{t}y(t),\label{eq:pico}\\
  y_{m}=\frac1T\int_{0}^{T}y(t)\,\dif t,\qquad
  y_{\mathrm{rms}}=\biggl(\frac1T\int_{0}^{T}y^{2}(t)\,\dif t\biggr)^{1/2}.\label{eq:rms}
\end{gather}
\end{definicion}

Para una señal simétrica de media nula, $y_{pp}=2y_{p}$, y el valor eficaz vale $y_p/\sqrt2$
(sinusoide), $y_p$ (onda cuadrada) o $y_p/\sqrt3$ (onda triangular). El osciloscopio mide con
comodidad $y_{pp}$; por eso, en lo que sigue, las ganancias se compararán como cocientes de
amplitudes o de tensiones pico a pico, que coinciden para señales de la misma forma.

% -------------------------------------------------------------------------
\subsection{El amplificador como cuadripolo}\label{sec:cuadripolo}
% -------------------------------------------------------------------------
Cualquier amplificador lineal, visto desde sus dos puertos (entrada y salida, con un terminal
común a tierra), se describe con el modelo de la \cref{fig:cuadripolo}: la entrada es una
impedancia $\Zin$ y la salida, un equivalente de Thévenin formado por una fuente controlada
$A_{vo}v_i$ y una impedancia $\Zout$~\cite{TEORIA-EA,TEORIA-ADICIONAL,SEDRA}.

% Figura: modelo de cuadripolo de un amplificador de tensión, con fuente y carga.
\begin{figure}[ht]
\centering
\begin{circuitikz}[circuito, scale=0.95, transform shape]
  % Recuadro del amplificador
  \fill[azulfondo] (2.8,-0.45) rectangle (7.3,2.65);
  \draw[azulclaro, line width=0.5pt, densely dashed] (2.8,-0.45) rectangle (7.3,2.65);
  \node[nota, anchor=south] at (5.05,2.65) {amplificador};
  % Fuente y resistencia interna
  \draw (0,0) to[sV, l=$v_s$] (0,2) to[R, l=$R_s$] (2.2,2)
        to[short, o-, i=$i_i$] (3.5,2) to[R, l=$\Rin$] (3.5,0) -- (2.2,0) to[short, -o] (2.2,0) -- (0,0);
  \draw (2.2,2) to[open, v^=$v_i$] (2.2,0);
  % Salida: equivalente de Thévenin
  \draw (3.5,0) -- (5.3,0) to[cV, l_=$A_{vo}v_i$] (5.3,2) to[R, l=$\Rout$] (6.9,2)
        to[short, -o, i=$i_o$] (7.9,2) -- (9.1,2) to[R, l=$R_L$] (9.1,0) -- (7.9,0)
        to[short, o-] (7.9,0) -- (5.3,0);
  \draw (7.9,2) to[open, v^=$v_o$] (7.9,0);
  \draw (4.4,0) node[ground] {};
  % Impedancias vistas desde cada puerto
  \draw[->, azulacero, line width=0.6pt] (1.6,2.85) -- node[above, font=\scriptsize, text=azulacero] {$\Zin$} (2.6,2.85);
  \draw[->, azulacero, line width=0.6pt] (8.5,2.85) -- node[above, font=\scriptsize, text=azulacero] {$\Zout$} (7.5,2.85);
\end{circuitikz}
\caption{Modelo de un amplificador de tensión como cuadripolo~\cite{TEORIA-EA,SEDRA}. La fuente
$v_s$ tiene resistencia interna $R_s$ (en el generador del laboratorio, $50\,\Omega$) y el
amplificador alimenta una carga $R_L$. $\Zin$ es la impedancia que ve la fuente y $\Zout$ la que ve
la carga.}
\label{fig:cuadripolo}
\end{figure}


\begin{definicion}[Parámetros de un amplificador]\label{def:parametros}
Con las magnitudes de la \cref{fig:cuadripolo},
\begin{equation}\label{eq:parametros}
  \Zin=\frac{v_i}{i_i},\qquad
  A_V=\frac{v_o}{v_i},\qquad
  A_I=\frac{i_o}{i_i},\qquad
  A_W=\frac{v_o\,i_o}{v_i\,i_i}=A_V\,A_I ,
\end{equation}
y $\Zout$ es la impedancia de Thévenin vista desde los terminales de salida con la fuente de
entrada anulada ($v_s=0$).
$A_{vo}$ es la ganancia en tensión \emph{en circuito abierto} ($R_L\to\infty$).
\end{definicion}

\begin{proposicion}[Efecto de carga]\label{prop:carga}
Si el amplificador se excita con una fuente de resistencia interna $R_s$ y se carga con $R_L$, la
ganancia efectiva desde la fuente es
\begin{equation}\label{eq:carga}
  \frac{v_o}{v_s}=\underbrace{\frac{\Rin}{\Rin+R_s}}_{\text{divisor de entrada}}\;A_{vo}\;
                  \underbrace{\frac{R_L}{R_L+\Rout}}_{\text{divisor de salida}} .
\end{equation}
\end{proposicion}
\begin{proof}
En la entrada, $R_s$ y $\Rin$ forman un divisor de tensión: $v_i=v_s\Rin/(\Rin+R_s)$. En la salida, la
fuente $A_{vo}v_i$ ve en serie $\Rout$ y $R_L$: $v_o=A_{vo}v_i\,R_L/(R_L+\Rout)$. Se multiplican ambas.
\end{proof}

\begin{corolario}[Amplificador de tensión ideal]\label{cor:ideal}
La ganancia no depende de la fuente ni de la carga si y solo si $\Rin\gg R_s$ y $\Rout\ll R_L$; en el
límite ideal, $\Rin\to\infty$ y $\Rout\to0$, y entonces $v_o/v_s=A_{vo}$.
\end{corolario}

Por eso hay que conocer \emph{las dos} impedancias además de una ganancia: la ganancia en potencia
$A_W=A_V^{2}R_{\mathrm{in}}/R_L$ (con $i_o=v_o/R_L$, $i_i=v_i/\Rin$) depende de su cociente.

\begin{definicion}[Ganancias en decibelios]\label{def:dB}
\begin{equation}\label{eq:dB}
  A_V(\mathrm{dB})=20\log_{10}\abs{A_V},\qquad
  A_I(\mathrm{dB})=20\log_{10}\abs{A_I},\qquad
  A_W(\mathrm{dB})=10\log_{10}A_W .
\end{equation}
\end{definicion}
El factor $20$ aparece porque la potencia es cuadrática en la tensión: así, $A_W(\mathrm{dB})=A_V(\mathrm{dB})$
cuando $\Rin=R_L$. Valores útiles: $\abs{A_V}=10\mapsto20\dB$, $100\mapsto40\dB$, $1\mapsto0\dB$ y
$1/\sqrt2\mapsto-3{,}01\dB$. Las ganancias pueden ser negativas (inversión de fase), de ahí el módulo; y,
salvo a baja frecuencia, son números complejos, de modo que hay que dar también su fase.
